% Source: https://artofproblemsolving.com/community/c3h2010801_2008_national_mathcounts_sprint_round_question_23
% Stable identity: aops-topic-2010801-q01
The right triangle has area $\dfrac{1}{2}\cdot3\cdot4=6$.

Let the second triangle be $PQR$ with $PQ=4$, $\{RP,RQ\}=\{x,y\}$, so $y=\max(RP,RQ)$. Its area is $6$, so the height from $R$ to line $PQ$ is $\dfrac{2\cdot6}{4}=3$: $R$ lies on a line parallel to $PQ$ at distance $3$.

\begin{asy}
usepackage("amsmath");
size(5.5cm);
defaultpen(fontsize(10pt));
pair P=(0,0), Q=(4,0);
pair R=(2,3), R2=(3.3,3);
draw((-1,3)--(5,3),dashed+gray);
draw(P--Q--R--cycle,linewidth(0.8));
draw(P--R2--Q,gray+dashed);
draw(R--(2,0),dotted);
draw((2,0)--(2,0.25)--(2.25,0.25)--(2.25,0));
dot(P);dot(Q);dot(R);dot(R2,gray);
label("$P$",P,SW); label("$Q$",Q,SE); label("$R$",R,N); label("$R'$",R2,N,gray);
label("$4$",(2,0),S);
label("$3$",(2,1.5),E);
label("$\sqrt{13}$",(P+R)/2,NW); label("$\sqrt{13}$",(Q+R)/2,NE);
label("\mbox{}",(min(currentpicture,user=true)+max(currentpicture,user=true))/2,black); // keeps the text after the figure black when LaTeX typesets the labels inline
\end{asy}

Put $P=(0,0)$, $Q=(4,0)$ and $R=(h,3)$. Then
\[RP^2=h^2+9,\qquad RQ^2=(4-h)^2+9.\]
Since $|h|+|4-h|\ge4$, at least one of $|h|$ and $|4-h|$ is at least $2$, so
\[y^2=\max\left(h^2,(4-h)^2\right)+9\ge4+9=13.\]
Equality holds for $h=2$, where $x=y=\sqrt{13}$ (and $\sqrt{13}$, $\sqrt{13}$, $4$ do form a triangle). A vertex such as $R'$ off the perpendicular bisector makes one side longer. Hence the smallest possible value of $y$ is $\boxed{\sqrt{13}}$.
